\section{Introduction and Objectives}
\label{sec:introduction}

Stage~1~\cite{stage1} built the data layer of our search engine over Project Gutenberg: a datalake with the \code{batch} layout
and an inverted index in MongoDB, chosen after benchmarking three languages and several structures. Stage~2~\cite{assignment} turns that
data layer into cooperating services. Each service runs as its own process, they share explicit contracts, and an
explicit mechanism coordinates them. Everything runs on one machine, and replicas and threads are part of the design
from the start. The goal is a system whose limits we can measure on that machine and whose design we can carry to the
Stage~3 cluster.

We set six objectives and met all of them:

\begin{enumerate}
\item \textbf{Four services and one module of contracts.} \code{crawler} (Ingestion), \code{indexer} (Indexing),
  \code{query} (Search) and \code{control} (seeding and reconciliation) are written in Java~25 with Javalin and
  Jackson. They share \code{common}, which is the only place where the code names a collection, a message, a state or a
  variable. MongoDB is the source of truth, the datalake is a Docker volume, and nginx is the only entry point.
\item \textbf{An explicit, failure-proof coordination mechanism.} Work travels as messages on Hazelcast queues, MongoDB
  decides who owns each job, and \code{control} periodically repairs whatever a crash leaves behind. A book becomes
  searchable only through a published batch, and it never appears in two published batches of the same generation.
\item \textbf{Documented and validated APIs.} The front-end, operational, health and internal interfaces are kept as
  separate inventories in \code{api/openapi.yaml} (OpenAPI~3.1, version~1.11). The non-HTTP contracts are in
  \code{SPEC.md}. An automated checker validates every response against the specification in continuous integration,
  in every coordination mode.
\item \textbf{Recovery from every failure we could cause.} The failure matrix has 14 rows, each with its own script. All
  165 recovery runs of the official campaign recover, and every row gives the same books and search answers with 1, 2
  and 3 replicas and in every mode where it applies.
\item \textbf{Alternatives measured, not only argued.} The same system can coordinate in four other ways behind one
  variable: ActiveMQ with and without persistence, polling of shared state, and classical orchestration. One campaign
  on one machine compared all five against each other.
\item \textbf{Benchmarks that can be traced.} JMH microbenchmarks and nine integration experiments ran with five
  interleaved repetitions and 95\,\% confidence intervals. Every figure in this report names the run it comes from,
  and that run keeps its command, commit, configuration and raw data in the repository (Appendix~\ref{app:runs}).
\end{enumerate}

The campaign answers the questions of the assignment. Indexing limits the pipeline: one indexer publishes 28.5
books/s, while one crawler downloads about 500. One \code{query} replica saturates at about 4{,}400 requests/s.
Coordination itself is cheap in every mode, below 17\,\% of one core. The default mode, \mode{hazelcast}, remains the
winner under the decision rule we fixed before the campaign (Section~\ref{sec:winner}).

Table~\ref{tab:team} shows how the team split the work. Each member owned one piece and reviewed another member's pull
requests, and changes to the shared contracts required all four. The work advanced through seven phases, F0 to F6,
and each phase had a gate: the walking skeleton had to show a new book in \code{/search}, and F3 required the whole
failure matrix to pass. The history of \code{dev} shows the contribution of each member.

\begin{table}[ht]
\centering
\small
\caption{Pieces owned by each member. Each owner wrote the code, scripts and runs of their piece.}
\label{tab:team}
\begin{tabularx}{\textwidth}{@{}lY@{}}
\toprule
Member & Piece \\
\midrule
Hugo González Fernández (A) & \code{crawler} (Ingestion), Docker Compose and the shared \code{Dockerfile}; \code{ingest-throughput} \\
Máximo Prosperi Reynoso (B) & \code{indexer} (Indexing and batch publication); \code{index-throughput}; the four alternative coordination modes, their experiments and the campaign script \\
Abel Castro Villar (C) & \code{query} (Search), nginx, \code{api/openapi.yaml} and the \code{checks} module; the load generator and \code{search-latency} \\
Javier Martín Rodríguez (D) & \code{control} (seeding and reconciliation), the failure matrix and its 14 recovery scripts, result processing in \code{bench}; \code{coordination-cost}, \code{resource-usage}, and the official campaign on his machine \\
\bottomrule
\end{tabularx}
\end{table}
